
\section*{\underline{EXPOSÉ X}\\[3pt] \underline{\underline{FORMALISME DES INTERSECTIONS SUR LES SCHÉMAS ALGÉBRIQUES PROPRES}}}
\begin{center}
\underline{par O. Jussila}\\
(avec un appendice par A. Grothendieck)
\end{center}

\section*{1. Compatibilité des filtrations avec la loi de composition}
@@M0@@

\subsection*{1.1.}
Considérons d'abord un schéma noethérien @@M1@@ et le groupe @@M2@@ des classes de @@M3@@-modules cohérents. Pour tout @@M4@@-module cohérent @@M5@@ nous désignerons par @@M6@@ sa classe dans @@M7@@.

\textbf{Définition 1.1.1.} \uline{Pour tout entier} @@M8@@ \uline{on désigne par} @@M9@@ \uline{le sous-groupe de} @@M10@@ \uline{engendré par les classes} @@M11@@ \uline{des} @@M12@@\uline{-modules cohérents} @@M13@@ \uline{tels que}
@@M14@@
On définit ainsi une filtration croissante @@M15@@ de @@M16@@ telle que @@M17@@ pour @@M18@@.

\textbf{Proposition 1.1.2.} \uline{Soient} @@M19@@ \uline{un schéma noethérien,} @@M20@@ \uline{un} @@M21@@\uline{-module cohérent tel que} @@M22@@\uline{, et} @@M23@@ (@@M24@@\uline{) le sous-ensemble} @@M25@@ \uline{de} @@M26@@\uline{. Pour tout} @@M27@@ \uline{on désigne par} @@M28@@ \uline{le sous-schéma intègre fermé associé à} @@M29@@\uline{. Alors la longueur du} @@M30@@\uline{-module} @@M31@@ \uline{est finie pour tout} @@M32@@ \uline{et non-nulle pour un nombre fini des} @@M33@@ \uline{seulement, et on a}
